%==============================================================================
%  Mid-Semester Examination : Scale Invariant Feature Transform
%  Report : NDSS-SIFT, a Gaussian-free alternative to SIFT
%
%  Compiles with pdflatex (Overleaf: default settings, no bibtex needed).
%  Put fig1_original.png ... fig8_invariance.png next to this file; they are
%  written automatically by the accompanying notebook / script.
%==============================================================================
\documentclass[11pt,a4paper]{article}

\usepackage[margin=2.4cm]{geometry}
\usepackage{amsmath,amssymb}
\usepackage{graphicx}
\usepackage{booktabs}
\usepackage{caption}
\usepackage{xcolor}
\usepackage{algorithm}
\usepackage{algpseudocode}
\usepackage[hidelinks]{hyperref}

\captionsetup{font=small,labelfont=bf}
\setlength{\parskip}{0.35em}

%------------------------------------------------------------------------------
%  RESULT NUMBERS -- update these SIX lines after the final run, nothing else.
%  (values below are from the run reported in the notebook output)
%------------------------------------------------------------------------------
\newcommand{\Nraw}{2507}          % raw scale-space extrema
\newcommand{\Nfilt}{1150}         % after weak / edge rejection
\newcommand{\Nloc}{1028}          % after sub-pixel localisation + de-duplication
\newcommand{\Nori}{1308}          % final oriented keypoints
\newcommand{\Truntime}{16.3}      % total runtime in seconds
\newcommand{\Kvals}{0.0319,\ 0.0090,\ 0.0092,\ 0.0140}  % per-octave contrast k

\title{\bfseries NDSS-SIFT: A Gaussian-Free Scale Invariant Feature Transform\\[2mm]
\large Nonlinear-Diffusion Scale Space with a Determinant-of-Hessian Response}
\author{\textit{Name} \\ \textit{Roll number} \\ \textit{Course / Instructor}}
\date{\today}

\begin{document}
\maketitle

\begin{abstract}
\noindent
This report presents \textbf{NDSS-SIFT}, an alternative to the Scale Invariant
Feature Transform in which \emph{every} Gaussian kernel of the original
algorithm has been removed, addressing part~(a) of the question. The starting
observation is that Gaussian smoothing is not a primitive operation but the
solution of the linear heat equation; the scale space is therefore rebuilt from
a \emph{nonlinear} Perona--Malik diffusion, which admits no convolution kernel
of any kind and hence no Gaussian. The Difference of Gaussians is replaced by
the scale-normalised determinant of the Hessian, the Gaussian pre-blur before
octave decimation is dropped entirely, and the Gaussian weighting windows of
the orientation histogram and of the $128$-D descriptor are replaced by
triangular (Bartlett) windows. All derivatives are central finite differences,
so not even the binomial $[1\ 2\ 1]$ factor hidden inside the Sobel operator
survives. The keypoint-selection machinery of SIFT --- $3\!\times\!3\!\times\!3$
extrema detection, contrast and principal-curvature rejection, quadratic
sub-pixel localisation, the $80\,\%$ orientation rule and the
$4\!\times\!4\!\times\!8$ descriptor --- is retained unchanged, so the method
remains recognisably SIFT. On the mandatory $512\times512$ Cameraman image the
complete pipeline, including an invariance study, runs in
$\approx\Truntime$\,s in Google Colab, well inside the three-minute budget.
\end{abstract}

%==============================================================================
\section{Problem statement}
%==============================================================================

SIFT~\cite{lowe2004} is built almost entirely out of Gaussian kernels. The task
is to design and implement an alternative that does not use them, with full
credit for eliminating \emph{all} of them rather than only replacing the
Difference of Gaussians (DoG) by a Laplacian pyramid. Table~\ref{tab:where}
lists where the Gaussians actually occur; any honest solution must deal with
all five sites, not just the DoG.

\begin{table}[ht]
\centering
\small
\begin{tabular}{@{}clp{8.2cm}@{}}
\toprule
\# & Site in classical SIFT & Role played by the Gaussian \\
\midrule
1 & Scale space $L(\cdot,\sigma)=G_\sigma * I$ & generates the family of
    increasingly coarse images \\
2 & Difference of Gaussians & band-pass response whose extrema are the keypoints \\
3 & Pre-blur before octave down-sampling & anti-aliasing before decimation by 2 \\
4 & Orientation histogram window & spatial weighting of gradient votes \\
5 & Descriptor window & spatial weighting of the $4\times4$ cell votes \\
\bottomrule
\end{tabular}
\caption{The five places where classical SIFT uses a Gaussian kernel.}
\label{tab:where}
\end{table}

%==============================================================================
\section{Key idea: a Gaussian is not a primitive}
%==============================================================================

The observation that makes a complete elimination possible is that Gaussian
smoothing \emph{is} linear diffusion. If $L$ evolves under the heat equation
\begin{equation}
\frac{\partial L}{\partial t}=\nabla^{2}L ,
\qquad L(\cdot,0)=I ,
\label{eq:heat}
\end{equation}
then the solution at time $t$ is exactly $L(\cdot,t)=G_{\sigma}*I$ with
$\sigma=\sqrt{2t}$. In other words the Gaussian is merely the Green's function
of~\eqref{eq:heat}; SIFT's scale space is a \emph{linear diffusion} scale space.

This reframing turns "remove the Gaussian" into "change the diffusion". We
therefore make the diffusion \emph{nonlinear}, in the sense of Perona and
Malik~\cite{perona1990}:
\begin{equation}
\boxed{\;
\frac{\partial L}{\partial t}
=\operatorname{div}\!\bigl(g(|\nabla L|)\,\nabla L\bigr),
\qquad
g(z)=\frac{1}{1+z^{2}/k^{2}} \; }
\label{eq:pm}
\end{equation}
Because the conductivity $g$ depends on the image being evolved, the evolution
operator is nonlinear and \emph{not shift invariant}. A nonlinear,
non-shift-invariant operator has no convolution representation at all: there is
no kernel to remove, Gaussian or otherwise. This is a strictly stronger
statement than "the Gaussian was swapped for another kernel".

The behaviour of~\eqref{eq:pm} is also better suited to a keypoint detector
than Gaussian blurring. Where $|\nabla L|\ll k$ the conductivity is
$g\approx1$ and homogeneous regions are smoothed; where $|\nabla L|\gg k$ we
get $g\approx0$ and the flow refuses to smooth \emph{across} object
boundaries. Structures therefore stay where they are as the scale increases,
instead of drifting and dissolving as they do under a Gaussian. The same
argument underlies the KAZE detector~\cite{kaze2012}; note, however, that KAZE
still computes a Gaussian-smoothed image to estimate its contrast factor and to
build its derivatives, whereas the present method does not.

\paragraph{The exponential conductivity is deliberately avoided.}
Perona and Malik also propose $g(z)=\exp(-z^{2}/k^{2})$. That function is a
Gaussian \emph{profile} evaluated on gradient magnitudes. Although it is not a
convolution kernel, the rational form in~\eqref{eq:pm} is used instead so that
no Gaussian-shaped expression appears anywhere in the implementation.

%==============================================================================
\section{The proposed algorithm}
%==============================================================================

\subsection{Nonlinear-diffusion scale space}
\label{sec:ss}

\paragraph{Scale ladder.} SIFT's octave structure is retained. With $S$ usable
levels per octave we build $S+3$ levels at octave-local scales
\begin{equation}
\sigma_i=\sigma_0\,2^{\,i/S},\qquad i=0,\dots,S+2 ,
\end{equation}
and reach level $i$ by integrating~\eqref{eq:pm} to the evolution time
\begin{equation}
t_i=\tfrac12\sigma_i^{2} .
\label{eq:tsigma}
\end{equation}
With $\sigma_0=1.6$ and $S=3$ this gives
$\sigma_i\in\{1.600,\,2.016,\,2.540,\,3.200,\,4.032,\,5.080\}$ and
$t_i\in\{1.28,\,2.03,\,3.23,\,5.12,\,8.13,\,12.90\}$. Levels are produced
incrementally, level $i$ being obtained from level $i-1$ by evolving for
$\Delta t_i=t_i-t_{i-1}$.

Relation~\eqref{eq:tsigma} is the exact time--scale correspondence of
\emph{linear} diffusion and is used here purely as a \emph{parameterisation}:
$\sigma_i$ is a label meaning ``the scale linear diffusion would have reached in
the same time'', chosen so that the parameters are directly comparable with
Lowe's. No Gaussian is ever evaluated; the symbol $\sigma$ could be eliminated
from the implementation without changing a single number.

\paragraph{Contrast factor.} The threshold $k$ separating "texture" from
"edge" is estimated automatically as the $70^{\text{th}}$ percentile of the
gradient-magnitude distribution of the octave's base image. It is re-estimated
once per octave, because decimation changes the gradient statistics. The values
obtained on Cameraman are $k=\Kvals$ for octaves $0$--$3$.

\paragraph{Discretisation.} Equation~\eqref{eq:pm} is integrated with an
explicit scheme on an \emph{isotropic} nine-point stencil,
\begin{equation}
L^{n+1}_{c}=L^{n}_{c}
+\Delta t\sum_{p\in\mathcal{N}_8(c)} w_p\,
\frac{g_p+g_c}{2}\,\bigl(L^{n}_{p}-L^{n}_{c}\bigr),
\qquad
w_p=\begin{cases}2/3 & \text{axial}\\[2pt] 1/6 & \text{diagonal}\end{cases}
\label{eq:scheme}
\end{equation}
The half-pixel conductivities $\tfrac12(g_p+g_c)$ make the scheme
conservative. For $g\equiv1$ the stencil~\eqref{eq:scheme} reduces to the
classical isotropic Laplacian $\frac16\left[\begin{smallmatrix}1&4&1\\
4&-20&4\\ 1&4&1\end{smallmatrix}\right]$, whose leading anisotropy error
cancels; the naive four-neighbour Laplacian is anisotropic on the pixel grid
and measurably degrades rotation invariance, since the scale space is what
decides \emph{where} keypoints appear. The weights sum to $10/3$, so the
scheme is stable for $\Delta t\le 3/10$; we use $\Delta t\le0.28$ and split
each increment into $n=\lceil\Delta t_i/0.28\rceil$ equal steps.

\paragraph{Octaves.} At the end of an octave, level $S$ (which carries
$\sigma=2\sigma_0$) is down-sampled by plain decimation, $L[::2,::2]$, with
\emph{no} filter of any kind. Decimation is grid-exact --- pixel $(i,j)$ of the
new octave is pixel $(2i,2j)$ of the old one --- so no half-pixel correction is
needed. Four octaves are built ($512,256,128,64$). The first level of octave
$o>0$ already carries $\sigma_0$ in the new coordinate system, so its evolution
starts from $t_0$ rather than from $0$.

\subsection{Response function: scale-normalised determinant of the Hessian}

Lowe uses the DoG because
$G_{k\sigma}-G_{\sigma}\approx(k-1)\sigma^{2}\nabla^{2}G$: the DoG is an
inexpensive approximation of the scale-normalised Laplacian. Since we no longer
have Gaussians to subtract, the differential operator is computed directly.
For every level we form the Hessian by finite differences and take
\begin{equation}
R(x,y,\sigma_i)=\sigma_i^{4}\,\det\mathbf{H}
=\sigma_i^{4}\bigl(L_{xx}L_{yy}-L_{xy}^{2}\bigr).
\label{eq:response}
\end{equation}
Three properties justify this choice.
\begin{enumerate}\itemsep2pt
\item \textbf{Scale selection.} Under a change of scale the entries of
$\mathbf{H}$ scale as $\sigma^{-2}$, so $\det\mathbf{H}$ scales as
$\sigma^{-4}$; the factor $\sigma^{4}$ makes $R$ scale invariant in the sense
of Lindeberg~\cite{lindeberg1998}, and a blob of characteristic size $\sigma$
produces a genuine maximum \emph{along the scale axis}. This is what allows
each keypoint to carry its own scale.
\item \textbf{Automatic edge suppression.} Along a straight edge one principal
curvature vanishes, so $\det\mathbf{H}\le0$ and the point cannot be a positive
maximum. The DoG, being a signed Laplacian-like operator, responds strongly to
edges and relies entirely on the subsequent curvature test to remove them.
\item \textbf{No kernel required.} $R$ is computed from three finite-difference
images, which is consistent with the Gaussian-free constraint.
\end{enumerate}

\subsection{Derivative operators}

All derivatives use central differences on the replicate-padded image:
\begin{align}
L_x &= \tfrac12\,[-1,0,1]*L,
\qquad L_{xx} = [1,-2,1]*L, \\[2pt]
L_{xy} &= \tfrac14
\begin{bmatrix}1&0&-1\\ 0&0&0\\ -1&0&1\end{bmatrix}*L .
\end{align}
The Sobel operator is deliberately \emph{not} used: it factorises as
$[-1,0,1]\otimes[1,2,1]$, and the binomial factor $[1,2,1]$ is the standard
\emph{discrete Gaussian} approximation. Using Sobel would reintroduce a
Gaussian into every derivative in the pipeline.

\subsection{Scale-space extrema (Figure~\ref{fig:raw})}

A sample is kept if it is a \emph{strict} maximum of $R$ among its $26$
neighbours in the $3\times3\times3$ neighbourhood spanning $(x,y,\sigma)$,
excluding the first and last level of each octave and a border of $6$ pixels.
Excluding the centre from the comparison footprint makes the test strict, so a
flat ridge of the response no longer yields a cluster of duplicate detections.
A numerical floor of $10^{-6}$ removes extrema caused by round-off in
homogeneous regions.

\subsection{Rejection of weak and unstable detections (Figure~\ref{fig:filt})}

Two tests, both taken from SIFT:
\begin{itemize}\itemsep2pt
\item \textbf{Contrast:} $R\ge\tau$, with $\tau=10^{-3}$ on an image scaled to
$[0,1]$.
\item \textbf{Principal curvature:} an unstable keypoint lies on a ridge, where
the ratio $r$ of the principal curvatures is large. Using
$\operatorname{tr}^2\mathbf{H}/\det\mathbf{H}=(r+1)^2/r$, we require
\begin{equation}
\frac{\bigl(L_{xx}+L_{yy}\bigr)^{2}}{L_{xx}L_{yy}-L_{xy}^{2}}
<\frac{(r+1)^{2}}{r},\qquad r=10 .
\end{equation}
The quantity is independent of the $\sigma^4$ normalisation, so the test is
identical to Lowe's even though the underlying response is not.
\end{itemize}

\subsection{Sub-pixel and sub-scale localisation (Figure~\ref{fig:loc})}

A $3$-D quadratic is fitted to $R$ around each detection,
\begin{equation}
R(\mathbf{z}+\delta)\approx R+\nabla R^{\!\top}\delta
+\tfrac12\delta^{\!\top}\mathbf{H}_R\,\delta ,
\qquad
\hat\delta=-\mathbf{H}_R^{-1}\nabla R ,
\end{equation}
with $\mathbf{z}=(x,y,s)$ and all derivatives of $R$ taken by finite
differences on the response volume. If any component of $\hat\delta$ exceeds
$\tfrac12$ the sample is re-centred on the neighbouring one and the fit is
repeated (at most five times); failures to converge inside the grid are
discarded, as are points whose interpolated response
$R+\tfrac12\nabla R^{\!\top}\hat\delta$ falls below $\tau$. The refined
sub-level index gives a continuous scale,
$\sigma=\sigma_0 2^{(s+\hat\delta_s)/S}\,2^{o}$.

\paragraph{Duplicate suppression.} The same structure is frequently detected in
two adjacent octaves or at two adjacent levels. Such detections carry nearly
identical descriptors, which defeats the ratio test used at matching time, so
keypoints within $1.5$\,px whose scales differ by less than a factor $1.3$ are
merged, the strongest being kept.

\subsection{Orientation assignment (Figure~\ref{fig:ori})}

For each keypoint, gradients are read from the scale-space level nearest to its
refined scale and accumulated into a $36$-bin histogram of
$\theta=\arctan2(L_y,L_x)$. Each sample is weighted by its gradient magnitude
times a \textbf{triangular (Bartlett) window}
\begin{equation}
w(r)=\max\!\Bigl(0,\;1-\frac{r}{R}\Bigr),\qquad R=4\sigma ,
\end{equation}
which replaces SIFT's Gaussian window. The Gaussian there performs no
Gaussian-specific function: it is a soft radial roll-off ensuring that distant
pixels contribute less and that nothing appears or disappears abruptly at the
window edge. Any smooth, decaying, compactly supported window serves, and the
Bartlett window --- a first-order B-spline, the convolution of two box windows
--- is the simplest such choice.

Two refinements reduce quantisation error, which otherwise propagates directly
into every descriptor:
\begin{itemize}\itemsep2pt
\item each sample is split \emph{linearly} between its two neighbouring
angular bins rather than assigned to one, removing the $10^{\circ}$ bias of
hard binning;
\item the histogram is smoothed circularly with the triangular kernel
$[1\,2\,3\,2\,1]/9$. Note that the \emph{binomial} kernel $[1\,4\,6\,4\,1]/16$,
the usual choice here, is a discrete Gaussian and is therefore avoided.
\end{itemize}
Following SIFT, the peak is located to sub-bin accuracy by fitting a parabola
to the peak bin and its two neighbours, and every local peak reaching $80\,\%$
of the maximum spawns an additional keypoint with the same position and scale
but a different orientation.

\subsection{Descriptor}

The descriptor follows SIFT exactly, with the Gaussian envelope replaced. The
patch is sampled on the pixel grid, each offset being rotated by $-\theta$ so
that the descriptor is expressed in the keypoint's own frame:
\begin{equation}
\begin{pmatrix}u\\v\end{pmatrix}
=\frac{1}{3\sigma}
\begin{pmatrix}\cos\theta & \sin\theta\\ -\sin\theta & \cos\theta\end{pmatrix}
\begin{pmatrix}x-x_0\\ y-y_0\end{pmatrix},
\end{equation}
so that one descriptor cell spans $3\sigma$ pixels and the $4\times4$ grid has
half-width $h=6\sigma$. Gradient orientations are likewise referred to
$\theta$. Votes are distributed by \emph{trilinear} interpolation over the two
neighbouring cells in each spatial direction and the two neighbouring
orientation bins, giving a $4\times4\times8=128$ dimensional vector. The
spatial weighting is again triangular,
\begin{equation}
w(r)=\max\!\Bigl(0,\;1-\frac{r}{2.5\,h}\Bigr).
\end{equation}
The factor $2.5$ matters: a cone that reached zero at the patch boundary would
give the four corner cells almost no weight and silently reduce the descriptor
to about twelve effective cells, whereas SIFT's Gaussian envelope leaves them
roughly $30\,\%$ of the central weight. With $2.5h$ the triangular profile
tracks the Gaussian profile closely over the whole patch.

Finally the vector is $\ell_2$ normalised, clipped at $0.2$ to suppress the
effect of large gradient magnitudes, and renormalised --- Lowe's illumination
robustness step --- after which the \emph{RootSIFT}
transform~\cite{rootsift2012} is applied,
\begin{equation}
\mathbf{d}\;\leftarrow\;\sqrt{\mathbf{d}/\lVert\mathbf{d}\rVert_{1}} ,
\end{equation}
so that Euclidean distances between descriptors become Hellinger distances
between the underlying histograms. The result is already $\ell_2$ normalised,
so the matching stage is unchanged.

\subsection{Summary of the algorithm}

\begin{algorithm}[ht]
\caption{NDSS-SIFT}
\label{alg:main}
\begin{algorithmic}[1]
\State $I\gets$ Cameraman, converted to \texttt{float64} in $[0,1]$
\For{$o=0$ to $O-1$}
  \State $k\gets$ $70^{\text{th}}$ percentile of $|\nabla L_{\text{base}}|$
  \For{$i=0$ to $S+2$}
    \State evolve \eqref{eq:pm} by $\Delta t_i=\tfrac12(\sigma_i^2-\sigma_{i-1}^2)$
           using scheme \eqref{eq:scheme} with $\Delta t\le0.28$
    \State $R_i\gets\sigma_i^{4}\bigl(L_{xx}L_{yy}-L_{xy}^{2}\bigr)$
  \EndFor
  \State $L_{\text{base}}\gets L_{S}[::2,::2]$ \Comment{no anti-alias filter}
\EndFor
\State $\mathcal{K}\gets$ strict maxima of $R$ over the $26$ neighbours in $(x,y,\sigma)$
\State discard $R<\tau$ and $\operatorname{tr}^2\mathbf{H}/\det\mathbf{H}\ge(r{+}1)^2/r$
\State refine each keypoint by a $3$-D quadratic fit; drop non-convergent ones
\State merge duplicates ($1.5$\,px, scale ratio $1.3$)
\ForAll{keypoints}
  \State build the $36$-bin orientation histogram with a $4\sigma$ Bartlett window
  \State emit one keypoint per peak above $80\,\%$ of the maximum
  \State compute the $4\times4\times8$ descriptor, normalise, clip, RootSIFT
\EndFor
\end{algorithmic}
\end{algorithm}

%==============================================================================
\section{What was changed with respect to the SIFT taught in class}
%==============================================================================

Table~\ref{tab:changes} is the central comparison. Everything not listed there
is unchanged from Lowe's algorithm.

\begin{table}[ht]
\centering
\small
\begin{tabular}{@{}p{4.3cm}p{5.1cm}p{5.1cm}@{}}
\toprule
Stage & Classical SIFT & NDSS-SIFT \\
\midrule
Scale space &
$L=G_\sigma*I$, i.e.\ \emph{linear} diffusion &
Perona--Malik \emph{nonlinear} diffusion, isotropic $9$-point explicit scheme;
no convolution kernel exists \\
\addlinespace[2pt]
Band-pass response &
$\mathrm{DoG}=L(k\sigma)-L(\sigma)$ &
$\sigma^{4}\det\mathbf{H}$ from finite differences \\
\addlinespace[2pt]
Octave down-sampling &
Gaussian pre-blur, then decimate &
plain decimation $L[::2,::2]$, no filter \\
\addlinespace[2pt]
Derivatives &
usually Sobel ($[-1,0,1]\otimes[1,2,1]$) &
central differences $[-1,0,1]/2$, $[1,-2,1]$ \\
\addlinespace[2pt]
Orientation window &
Gaussian, $\sigma_w=1.5\sigma$ &
triangular $\max(0,1-r/4\sigma)$ \\
\addlinespace[2pt]
Histogram smoothing &
box/binomial &
triangular $[1\,2\,3\,2\,1]/9$, plus linear soft binning \\
\addlinespace[2pt]
Descriptor window &
Gaussian, $\sigma_w=h$ &
triangular $\max(0,1-r/2.5h)$ \\
\addlinespace[2pt]
Descriptor normalisation &
$\ell_2$, clip $0.2$, $\ell_2$ &
same, followed by RootSIFT \\
\addlinespace[2pt]
Duplicate handling &
none &
merge within $1.5$\,px and scale ratio $1.3$ \\
\midrule
\multicolumn{3}{@{}l@{}}{\emph{Retained unchanged:}
$3\!\times\!3\!\times\!3$ extrema search; contrast threshold; principal-curvature
test with $r=10$;}\\
\multicolumn{3}{@{}l@{}}{\ \ quadratic sub-pixel/sub-scale refinement;
$36$-bin orientation histogram with the $80\,\%$ rule; $4\!\times\!4\!\times\!8$}\\
\multicolumn{3}{@{}l@{}}{\ \ descriptor with trilinear interpolation; octave
structure with $S=3$, $\sigma_0=1.6$.}\\
\bottomrule
\end{tabular}
\caption{Differences between classical SIFT and the proposed method.}
\label{tab:changes}
\end{table}

\paragraph{Why these substitutions are legitimate.} The Gaussian plays two
genuinely different roles in SIFT. In the scale space it is \emph{structural}:
it is the unique kernel satisfying the scale-space axioms (linearity,
shift/rotation invariance, semi-group, no creation of new extrema), so it
cannot simply be swapped for another kernel --- the construction itself has to
change, which is why a nonlinear PDE is used. In the orientation and descriptor
windows it is merely a \emph{soft aperture}, and nothing in SIFT's derivation
depends on its Gaussian shape; a compactly supported triangular window
discharges the same duty.

%==============================================================================
\section{Gaussian-freeness: audit and objections}
%==============================================================================

Table~\ref{tab:audit} audits every operation in the implementation that touches
pixel values.

\begin{table}[ht]
\centering
\small
\begin{tabular}{@{}llc@{}}
\toprule
Operation & What is actually computed & Gaussian? \\
\midrule
first derivatives            & $[-1,0,1]/2$                      & no \\
second derivatives           & $[1,-2,1]$, $4$-point cross       & no \\
scale space                  & nonlinear PDE \eqref{eq:pm}       & \textbf{no kernel exists} \\
blob response                & $\sigma^4\det\mathbf{H}$          & no \\
octave step                  & decimation $[::2,::2]$            & no \\
non-maximum suppression      & morphological max (rank filter)   & no \\
histogram smoothing          & $[1\,2\,3\,2\,1]/9$ (Bartlett)    & no \\
orientation / descriptor window & $\max(0,1-r/R)$                & no \\
descriptor binning           & trilinear (linear weights)        & no \\
evaluation warps             & bilinear interpolation            & no \\
\bottomrule
\end{tabular}
\caption{Audit of every pixel-level operation. The implementation contains no
call to \texttt{exp}, no \texttt{gaussian\_filter}, no Sobel operator, and
never enables \texttt{anti\_aliasing} (which is Gaussian inside
\texttt{scikit-image}).}
\label{tab:audit}
\end{table}

Three objections deserve an explicit answer.

\paragraph{"The symbol $\sigma$ is still used."}
$\sigma$ labels evolution time through $t=\sigma^2/2$ and is retained only so
that the parameters are directly comparable with Lowe's. Nothing in the
implementation evaluates $\exp(-x^{2}/2\sigma^{2})$; driving the code by $t$
alone would produce bit-identical output.

\paragraph{"A triangular window is an approximate Gaussian."}
It is not. A Bartlett window is compactly supported and piecewise linear --- a
first-order B-spline, obtained by convolving two box windows. What \emph{would}
be a discrete Gaussian is the binomial family $[1\,2\,1]$, $[1\,4\,6\,4\,1]$,
which is precisely why those kernels are avoided throughout, including inside
the Sobel operator.

\paragraph{"For $|\nabla L|\ll k$ the PDE reduces to the heat equation."}
Locally, in perfectly homogeneous regions, the behaviour does approach linear
diffusion. Two points answer this. First, the operator is not shift invariant,
so no convolution kernel represents the evolution: there is nothing to
eliminate because nothing exists. Second, the measured contrast factors,
$k=\Kvals$, are small, so essentially every structure in Cameraman that
produces a keypoint lies in the strongly nonlinear regime; homogeneous sky
behaving diffusively is irrelevant, since it generates no keypoints.

%==============================================================================
\section{Implementation and complexity}
%==============================================================================

The implementation is pure \texttt{NumPy}; \texttt{SciPy} is used only for the
non-maximum-suppression filter and a $k$-d tree in the evaluation code, and
\texttt{scikit-image} only to load the image and to warp it in the invariance
study. Every stage is vectorised except the per-keypoint loops.

Let $N$ be the number of pixels of an octave. One explicit step
\eqref{eq:scheme} costs $\mathcal{O}(N)$; reaching level $i$ requires
$\lceil\Delta t_i/0.28\rceil$ steps, about $50$ steps per octave in total.
Since octave areas form a geometric series with ratio $1/4$, the whole scale
space costs $\mathcal{O}(N)$ with a small constant. Detection, filtering and
refinement are $\mathcal{O}(N)$ and $\mathcal{O}(|\mathcal{K}|)$; orientation
and descriptor cost $\mathcal{O}(|\mathcal{K}|\sigma^2)$.

\begin{table}[ht]
\centering
\small
\begin{tabular}{@{}llll@{}}
\toprule
Parameter & Value & Parameter & Value \\
\midrule
octaves $O$                 & $4$        & contrast threshold $\tau$      & $10^{-3}$ \\
levels per octave $S$       & $3$ ($S{+}3=6$ built) & curvature limit $r$ & $10$ \\
base scale $\sigma_0$       & $1.6$      & border margin                  & $6$\,px \\
contrast percentile         & $70$       & refinement iterations          & $5$ \\
time step $\Delta t$        & $\le0.28$  & duplicate tolerance            & $1.5$\,px, $1.3\times$ \\
stencil weights             & $2/3,\ 1/6$& orientation bins               & $36$ \\
response floor              & $10^{-6}$  & orientation window             & $4\sigma$ \\
descriptor grid             & $4\times4\times8$ & descriptor cell         & $3\sigma$ \\
descriptor envelope         & $2.5h$     & peak ratio                     & $0.8$ \\
\bottomrule
\end{tabular}
\caption{Complete parameter set. All values are defined in one configuration
block at the top of the notebook.}
\label{tab:params}
\end{table}

%==============================================================================
\section{Results}
%==============================================================================

All results are produced on the mandatory image,
\texttt{skimage.data.camera()}, $512\times512$, $8$-bit grayscale, converted to
\texttt{float64} in $[0,1]$. The complete notebook, including the invariance
study, executes in $\approx\Truntime$\,s from a fresh Colab runtime.

\begin{figure}[ht]
\centering
\includegraphics[width=0.46\textwidth]{fig1_original.png}
\caption{The mandatory $512\times512$ Cameraman image.}
\label{fig:orig}
\end{figure}

\begin{figure}[ht]
\centering
\includegraphics[width=\textwidth]{fig2_scale_space.png}
\caption{Multiscale responses of the proposed scale-space
construction. Top row: levels of the nonlinear-diffusion scale space. Bottom
row: the corresponding scale-normalised determinant-of-Hessian responses. The
scale $\sigma$ of every panel is printed above it, in pixels of the original
image; panels span roughly $\sigma=2$ to $\sigma=26$\,px across four octaves.
Note that region boundaries remain sharp as $\sigma$ grows --- the signature of
nonlinear diffusion, and the reason keypoints do not drift with scale.}
\label{fig:ss}
\end{figure}

\begin{figure}[ht]
\centering
\includegraphics[width=0.56\textwidth]{fig3_raw_extrema.png}
\caption{Raw scale-space extrema before any rejection or
refinement. \textbf{Total number of detected extrema: \Nraw.}}
\label{fig:raw}
\end{figure}

\begin{figure}[ht]
\centering
\includegraphics[width=0.56\textwidth]{fig4_filtered.png}
\caption{Keypoints surviving the contrast and
principal-curvature tests ($\Nfilt$ remain). Marker size is proportional to the
detected scale.}
\label{fig:filt}
\end{figure}

\begin{figure}[ht]
\centering
\includegraphics[width=0.56\textwidth]{fig5_localized.png}
\caption{Keypoints after spatial and scale refinement by the
$3$-D quadratic fit, followed by duplicate suppression ($\Nloc$ remain). Marker
size reflects the refined scale; the red dot marks the sub-pixel position.}
\label{fig:loc}
\end{figure}

\begin{figure}[ht]
\centering
\includegraphics[width=0.62\textwidth]{fig6_oriented.png}
\caption{Final oriented keypoints ($\Nori$ in total; the
count exceeds Figure~\ref{fig:loc} because a keypoint with several dominant
orientations is duplicated, as in SIFT). Each circle has radius $2\sigma$ and
each radial line shows the assigned orientation.}
\label{fig:ori}
\end{figure}

\begin{table}[ht]
\centering
\small
\begin{tabular}{@{}lrl@{}}
\toprule
Stage & Keypoints & Figure \\
\midrule
raw scale-space extrema                       & \Nraw  & \ref{fig:raw} \\
after weak-response and edge rejection        & \Nfilt & \ref{fig:filt} \\
after sub-pixel localisation and de-duplication & \Nloc  & \ref{fig:loc} \\
final oriented keypoints                      & \Nori  & \ref{fig:ori} \\
\bottomrule
\end{tabular}
\caption{Keypoint population at each stage of the pipeline.}
\label{tab:counts}
\end{table}

\subsection{Invariance study}

To verify that the detector really is scale and rotation invariant, the
pipeline is re-run on warped copies of Cameraman with a known ground-truth
transform. Two standard measures are reported. \emph{Repeatability} is the
fraction of keypoints of the common region that are re-detected within
$2.5$\,px and within a factor $1.5$ in scale; it measures the detector alone.
\emph{Matching precision} is the fraction of accepted descriptor matches that
are geometrically correct within $3$\,px, and the \emph{matching score} is the
number of correct matches divided by the number of available keypoints. Only
the common valid region is used: the black corners created by a rotation are
excluded, together with a $24$\,px margin so that a keypoint's support lies
entirely inside the valid area. Warping uses bilinear interpolation with
anti-aliasing disabled, since \texttt{scikit-image}'s anti-aliasing filter is
Gaussian.

Matching uses Lowe's ratio test at $0.8$ with two modifications. The second
nearest neighbour is required to lie at least $4$\,px away from the first,
because otherwise the same physical point detected twice acts as its own
second neighbour and a correct match is rejected; and matches are required to
be mutual nearest neighbours.

\begin{table}[ht]
\centering
\small
\begin{tabular}{@{}lrrrrrr@{}}
\toprule
Test & kp (A) & kp (B) & repeatability & matches & precision & m-score \\
\midrule
rotated $15^{\circ}$ & & & & & & \\
rotated $30^{\circ}$ & & & & & & \\
rotated $45^{\circ}$ & & & & & & \\
rotated $90^{\circ}$ & & & & & & \\
zoom $\times1.50$    & & & & & & \\
zoom $\times0.75$    & & & & & & \\
\bottomrule
\end{tabular}
\caption{Invariance study. \textcolor{red}{Fill in from the table printed by
the notebook.} The $90^{\circ}$ row is the sanity check: that rotation is exact
on the pixel grid, so repeatability there should approach $100\,\%$.}
\label{tab:bench}
\end{table}

\begin{figure}[ht]
\centering
\includegraphics[width=0.92\textwidth]{fig7_match_example.png}
\caption{\emph{Additional figure.} Correctly matched keypoints between
the original image and a $30^{\circ}$ rotated copy.}
\label{fig:match}
\end{figure}

\begin{figure}[ht]
\centering
\includegraphics[width=0.92\textwidth]{fig8_invariance.png}
\caption{\emph{Additional figure.} Repeatability, matching precision and
matching score as a function of rotation angle (left) and zoom factor (right).}
\label{fig:inv}
\end{figure}

%==============================================================================
\section{Discussion and limitations}
%==============================================================================

\paragraph{What nonlinear diffusion buys.} Because the flow does not smooth
across boundaries, the position of a structure is far more stable across scale
than under Gaussian blurring, where every feature drifts as it dissolves. This
is visible in Figure~\ref{fig:ss}: the silhouette of the cameraman is still
crisp at $\sigma\approx26$\,px.

\paragraph{What it costs.} Two limitations follow from the same property.
First, an edge-preserving flow leaves sharp edges intact, so decimation without
an anti-alias filter does alias along those edges; the implementation therefore
offers an optional $3\times3$ \emph{box} average before decimation (a box
kernel is not a Gaussian) which trades a little sharpness for less aliasing.
Second, because a strong structure never fully dissolves, its response can keep
growing with $\sigma$, which makes scale selection less decisive than under
Gaussian blurring --- the price paid for positional stability.

\paragraph{Scale range.} Unlike Lowe's implementation, the input image is not
up-sampled by a factor of two before processing, so structures finer than
$\sigma_0=1.6$\,px are not represented. This is why zooming \emph{out} is the
harder direction in Table~\ref{tab:bench}.

\paragraph{Cost of explicit integration.} The explicit scheme is limited to
$\Delta t\le0.3$, so reaching large scales requires many small steps. This is
comfortably fast at $512\times512$, but for large images an Additive Operator
Splitting or Fast Explicit Diffusion scheme~\cite{weickert1998} would remove
the step-size restriction without introducing any kernel.

%==============================================================================
\section{Conclusion}
%==============================================================================

NDSS-SIFT replaces every Gaussian kernel of SIFT. The scale space is rebuilt
from a nonlinear Perona--Malik diffusion, which has no convolution kernel at
all rather than a non-Gaussian one; the Difference of Gaussians becomes a
scale-normalised determinant of the Hessian computed from central differences;
the anti-alias pre-blur is removed; and the Gaussian weighting windows become
triangular. The selection, refinement, orientation and description stages of
SIFT are retained unchanged, so the method is a genuine alternative SIFT rather
than a different detector. The implementation produces all six mandatory
figures on the prescribed Cameraman image and executes end to end in
$\approx\Truntime$\,s in Google Colab.

%==============================================================================
\begin{thebibliography}{9}
\bibitem{lowe2004} D.~G. Lowe, ``Distinctive image features from
scale-invariant keypoints,'' \emph{IJCV}, 60(2):91--110, 2004.
\bibitem{perona1990} P.~Perona and J.~Malik, ``Scale-space and edge detection
using anisotropic diffusion,'' \emph{IEEE TPAMI}, 12(7):629--639, 1990.
\bibitem{weickert1998} J.~Weickert, \emph{Anisotropic Diffusion in Image
Processing}, Teubner, 1998.
\bibitem{lindeberg1998} T.~Lindeberg, ``Feature detection with automatic scale
selection,'' \emph{IJCV}, 30(2):79--116, 1998.
\bibitem{kaze2012} P.~F. Alcantarilla, A.~Bartoli and A.~J. Davison, ``KAZE
features,'' \emph{ECCV}, 2012.
\bibitem{rootsift2012} R.~Arandjelovi\'c and A.~Zisserman, ``Three things
everyone should know to improve object retrieval,'' \emph{CVPR}, 2012.
\end{thebibliography}

\end{document}
